% Part: applied-modal-logics
% Chapter: epistemic-logic
% Section: language-epistemic-logic
%READERNOTE{SOL6-A188: मोडल-विरहिततेची अट सर्व कर्त्यांसाठी केली. समूहज्ञानाचा संयोग फक्त सीमित समूहासाठी साधा formula abbreviation आहे हे स्पष्ट केले; अनंत समूहाचे स्वतंत्र अर्थनिर्धारण व सामायिक ज्ञानाचा भाषा-विस्तार वेगळा दिला. वैयक्तिक आणि सामायिक ज्ञान नेहमी काटेकोरपणे वेगळे असतात हा अर्थ टाळला.}

\documentclass[../../../include/open-logic-section]{subfiles}

\begin{document}

\olfileid{aml}{el}{lan}

\olsection{ज्ञानविषयक तर्कशास्त्राची भाषा}

\begin{defn}
$G$ हा कर्तृचिन्हांचा संच असू द्या. अनेक कर्त्यांच्या
ज्ञानविषयक तर्कशास्त्राच्या मूलभूत भाषेत पुढील गोष्टी
असतात:
\begin{enumerate}
  \tagitem{prvFalse}{!!{falsity} दर्शवणारा विधानीय स्थिरांक~$\lfalse$.}{}
  \tagitem{prvTrue}{!!{truth} दर्शवणारा विधानीय स्थिरांक~$\ltrue$.}{}
  \item !!{propositional variable}s चा !!{denumerable}s संच: $\Obj
    p_0$, $\Obj p_1$, $\Obj p_2$, \dots
  \item विधानीय संयोजके: \startycommalist
  \iftag{prvNot}{\ycomma $\lnot$ (नकार)}{}%
  \iftag{prvAnd}{\ycomma $\land$ (संयोग)}{}%
  \iftag{prvOr}{\ycomma $\lor$ (विकल्पयोग)}{}%
  \iftag{prvIf}{\ycomma $\lif$ (!!{conditional})}{}%
  \iftag{prvIff}{\ycomma $\liff$ (!!{biconditional})}{}.
  \item $a \in G$ असताना ज्ञान-संकारक $\Knows_a$.
\end{enumerate}
\end{defn}

आपल्या व्यवस्थेत आपल्याला एकाच कर्त्याच्या ज्ञानाचा
विचार करायचा असेल, तर संच~$G$ आणि स्वतंत्र कर्ते
यांचा निर्देश वगळता येतो. त्या वेळी आपल्याकडे फक्त
मूलभूत संकारक~$\Knows$ असतो.

\begin{defn}
ज्ञानविषयक भाषेतील \emph{!!^{formula}s} ची पुढीलप्रमाणे
  आगमनाने व्याख्या करतो:
\begin{enumerate}
\tagitem{prvFalse}{$\lfalse$ हे अण्विक !!{formula} आहे.}{}

\tagitem{prvTrue}{$\ltrue$ हे अण्विक !!{formula} आहे.}{}

\item प्रत्येक विधानीय चल~$\Obj p_i$ हे (अण्विक) !!{formula} आहे.

\tagitem{prvNot}{$!A$ हे !!a{formula} असेल, तर $\lnot !A$
  हेही !!a{formula} आहे.}{}

\tagitem{prvAnd}{$!A$ आणि $!B$ ही !!{formula}s असतील, तर $(!A \land
  !B)$ हे !!a{formula} आहे.}{}

\tagitem{prvOr}{$!A$ आणि $!B$ ही !!{formula}s असतील, तर $(!A \lor !B)$
  हे !!a{formula} आहे.}{}

\tagitem{prvIf}{$!A$ आणि $!B$ ही !!{formula}s असतील, तर $(!A \lif !B)$
  हे !!a{formula} आहे.}{}

\tagitem{prvIff}{$!A$ आणि $!B$ ही !!{formula}s असतील, तर $(!A \liff !B)$
  हे !!a{formula} आहे.}{}

\item $!A$ हे !!a{formula} असेल आणि $a \in G$ असेल,
  तर $\Knows_a !A$ हे !!a{formula} आहे.

\tagitem{limitClause}{याखेरीज दुसरे काहीही !!a{formula} नाही.}{}
\end{enumerate}
\end{defn}

!!a{formula}~$!A$ मध्ये कोणत्याही $a\in G$ साठी
$\Knows_a$ नसेल, तर त्याला \emph{मोडल-विरहित} म्हणतो.

\begin{defn}
  $\Knows$ हा संकारक वैयक्तिक ज्ञान दर्शवण्यासाठी
  आहे; $\EKnows$ हा ``सर्वांना माहीत आहे'' असा
  वाचला जाणारा संकारक समूहज्ञान दर्शवतो. रिकामेतर
  सीमित $G' \subseteq G$ साठी $\EKnows_{G'} !A$ हे
  पुढील सीमित संयोगाचे संक्षिप्त रूप म्हणून परिभाषित करतो:
  \[\bigwedge_{b \in G'} \Knows_b !A.\]
  रिकाम्या समूहासाठी हे सत्य असलेला रिकामा संयोग मानतो.
  $G'$ अनंत असेल, तर हा संयोग मूलभूत सीमित सूत्रांच्या भाषेतील
  संक्षेप नाही. अशा समूहासाठी स्वतंत्र $\EKnows_{G'}$
  संकारक जोडल्यास त्याची सत्यतेची अट अशी घ्यायची:
  $\mSat{M}{\EKnows_{G'} !A}[w]$ जर आणि फक्त जर
  प्रत्येक $b\in G'$ साठी $\mSat{M}{\Knows_b !A}[w]$.
\end{defn}

कर्त्यांच्या~$G$ समूहातील ज्ञानाची अधिक प्रबळ
संकल्पनाही विचारात घेता येते: \emph{सामायिक ज्ञान}.
यात त्या समूहातील कर्त्यांच्या प्रत्येक रिकामेतर सीमित
क्रमासाठी, एकाला माहीत आहे की दुसऱ्याला माहीत आहे की
आणखी एकाला माहीत आहे, अशा प्रकारचे ज्ञान आवश्यक असते;
अशी अट प्रत्येक सीमित लांबीसाठी घेतली जाते.
सामायिक ज्ञानातून समूहज्ञान मिळते, पण उलट दिशा सर्व
प्रतिमानांत खरी नसते. तरीही प्रत्येक समूह व चौकटीत
हा भेद काटेकोर असेलच असे नाही; उदाहरणार्थ, एकाच
कर्त्याचा परावर्ती संक्रमक संबंध असेल, तर त्याचे
वैयक्तिक आणि सामायिक ज्ञान जुळते.
``$!A$ हे~$G$ समूहात सामायिक ज्ञान आहे'' हे
दर्शवण्यासाठी आपण $\CKnows_G !A$ वापरू. सर्व सीमित
लांबींच्या अटींचा संयोग स्वतः मूलभूत भाषेतील सीमित
सूत्र नाही; म्हणून $\CKnows_G$ हा विस्तारित भाषेतील
स्वतंत्र संकारक मानतो. त्याचे संबंधी अर्थनिर्धारण
\olref[trw]{sec} मध्ये दिले आहे.

\end{document}
